\section{CoM 高さ参照の更新アルゴリズム}
\label{sec:comz_update_algorithm}

前節では，CoM の高さが高くなるほど支持脚を支点とした振り子運動が緩やかになり，転倒方向への傾きの進行が抑制されることを力学的に示した．本節では，この考察に基づき，床沈み込みを検知した際に CoM の鉛直方向目標値を能動的に増加させる歩行制御アルゴリズムを提案する．

\subsection{従来手法における CoM 高さ参照}
\label{subsec:comz_conventional}

従来の歩行制御手法では，CoM の鉛直方向目標値はあらかじめ生成された歩行計画軌道に追従する形で与えられる．時刻 $t$ における CoM の鉛直方向参照値を $z_{\mathrm{ref}}^{\mathrm{conv}}(t)$，歩行生成によって得られる計画軌道を $z_{\mathrm{plan}}(t)$ とすると，従来手法は次式で表される．
\begin{equation}
z_{\mathrm{ref}}^{\mathrm{conv}}(t)=z_{\mathrm{plan}}(t).
\label{eq:comz_conv}
\end{equation}
このように，従来手法では床環境の変化や床沈み込みが生じた場合であっても，CoM の鉛直方向参照は現在時刻の計画値に追従するのみであり，床沈み込みにより実環境側で CoM 高さが低下する状況に対して，参照値側が能動的に補償を行う仕組みを持たない．

\subsection{提案手法：沈み込み検知に基づく CoM 高さ参照の更新}
\label{subsec:comz_proposed}

提案手法では，床沈み込みが検知された瞬間に CoM の鉛直方向参照値を更新し，支持脚を支点とした振り子運動を緩やかにすることで転倒を抑制する．具体的には，床沈み込みを検知した時点で，CoM 高さ参照に対して鉛直方向速度および加速度に基づく先読み量を加算する．制御周期を $\Delta t$，CoM 鉛直方向の参照速度および参照加速度をそれぞれ $\dot{z}_{\mathrm{ref}}(t)$，$\ddot{z}_{\mathrm{ref}}(t)$ とし，床沈み込み検知を表すフラグを $s(t)\in\{0,1\}$ とすると，提案手法における CoM 高さ参照の更新則は次式で表される．
\begin{equation}
z_{\mathrm{ref}}^{\mathrm{prop}}(t+\Delta t)=z_{\mathrm{plan}}(t+\Delta t)+s(t)\left(\dot{z}_{\mathrm{ref}}(t)\Delta t+\frac{1}{2}\ddot{z}_{\mathrm{ref}}(t)\Delta t^2\right).
\label{eq:comz_prop}
\end{equation}
式 \eqref{eq:comz_prop} において，床沈み込みが検知されていない場合（$s(t)=0$）には従来手法と同様に計画軌道に追従し，床沈み込みが検知された場合（$s(t)=1$）にのみ CoM 高さ参照が能動的に増加する．この操作により，床沈み込みが生じる局面において CoM 高さを一時的に高く保つことができ，支持脚を支点とした回転運動の角加速度を低減することで，DS から SS への遷移時に発生する過大なロール方向支持脚トルクの抑制が期待される．

\subsection{従来手法と提案手法の比較}
\label{subsec:comz_compare}

従来手法と提案手法における CoM 高さ参照の更新方法の違いを Table~\ref{tab:comz_compare} に示す．従来手法は計画軌道追従のみであるのに対し，提案手法は床沈み込み検知をトリガとして CoM 高さ参照を先読み的に更新する点に特徴がある．

\begin{table}[t]
  \centering
  \caption{CoM 高さ参照の更新則の比較}
  \label{tab:comz_compare}
  \begin{tabular}{l l}
    \hline
    手法 & CoM 高さ参照の更新方法 \\
    \hline
    従来手法 &
    $z_{\mathrm{ref}}^{\mathrm{conv}}(t)=z_{\mathrm{plan}}(t)$ \\
    提案手法 &
    $z_{\mathrm{ref}}^{\mathrm{prop}}(t+\Delta t)=z_{\mathrm{plan}}(t+\Delta t)+s(t)\left(\dot{z}_{\mathrm{ref}}(t)\Delta t+\frac{1}{2}\ddot{z}_{\mathrm{ref}}(t)\Delta t^2\right)$ \\
    \hline
  \end{tabular}
\end{table}
